% Gauss, Werke X/1 (1917), native99–102, printed92–95.
\SourceLeaf{099}\NoSourceHead
\SourceGroup{V}
\HistoricTitle{[ZUR THEORIE DER FORMEN.]}
\HistoricCitation{[Ein Zettel in Ed3, Kapsel 44.]}
\NumberedItem{1}
\noindent\hspace*{\parindent}Wir bezeichnen mit $p$ eine als Modulus anzusehende Primzahl; mit
$x,y,z\,.\,.\,.$ unbestimmte Zahlen ohne Einschränkung (oder lieber $0,1,2,.\,.\,.,
p-1$); mit $t,u,v\,.\,.\,.$ unbestimmte Zahlen durch $p$ nicht theilbar (oder lieber
$1,2,3,.\,.\,.,p-1$); durch $R$ einen bestimmten quadratischen Rest, durch $N$
einen bestimmten Nichtrest. Von den doppelten Ansätzen bezieht sich der
erste auf den Fall $p\equiv1$, der andere auf den Fall $p\equiv3$ (mod. 4).
\begin{center}Anzahl der Auflösungen der Congruenzen.\end{center}
\[
\begin{array}{lr|l}
xx&\equiv0&1\\
xx+yy&\equiv0&2p-1,\ 1\\
xx-yy&\equiv0&2p-1\\
xx+yy-zz&\equiv0&pp\\
xx+yy+zz&\equiv0&pp\\
tt+uu+R&\equiv0&p-5,\ p+1\\
tt+uu+N&\equiv0&p-1,\ p-3\\
tt+uu&\equiv0&2p-2,\ 0\\
\llap{\smash{\raisebox{-.5\normalbaselineskip}{$\left\{\vphantom{\begin{array}{c}x\\x\end{array}}\right.$}}}xx+yy+R&\equiv0&p-1,\ p+1\\
xx+yy+N&\equiv0&p-1,\ p+1\\
tt-uu+R&\equiv0&p-5,\ p-3\\
tt-uu+N&\equiv0&p-1,\ p-3\\
tt-uu&\equiv0&2p-2
\end{array}
\]
\clearpage
\SourceLeaf{100}\ArithmeticHead{93}{R}{ZUR THEORIE DER FORMEN.}
\[
\begin{array}{lr|l}
xx-yy+R&\equiv0&p-1\\
xx-yy+N&\equiv0&p-1\\
xx+x'x'+x''x''+x'''x'''&\equiv0&p^3+pp-p\\
xx+yy+zz+R&\equiv0&pp+p\\
xx+yy+zz+N&\equiv0&pp-p\\
xx+x'x'+x''x''-x'''x'''&\equiv0&p^3+pp-p,\ p^3-pp+p\\
xx+yy-zz+R&\equiv0&pp+p,\ pp-p\\
xx+yy-zz+N&\equiv0&pp-p,\ pp+p.
\end{array}
\]
\NumberedItem{2}
\noindent\hspace*{\parindent}Die Schlüsse von \textsc{Lagrange}, \textit{Mem. de l’Ac. de Berlin} 1770, p. 126 werden
durch Heranziehung der complexen Grössen ungemein vereinfacht.

Bezeichnen wir durch $p',q'$ die Adjuncten der complexen Zahlen $p,q$,
so haben wir die ursprüngliche Gleichung
\[
Aa=pp'+qq',
\]
\noindent die wir Kürze halber auch so schreiben
\[
Aa=(p,q).
\]
\noindent Es ist Voraussetzung, dass $p$ und $a$ inter se primi sind; man mache also
\[
m\equiv\frac{q}{p}\quad(\mathrm{mod.}\ a)
\]
\noindent und setze
\[
1+mm'=ab.
\]
\noindent Dann ist in ganzen Zahlen
\[
Ab=\left(\frac{p+mq}{a},\ \frac{q-pm}{a}\right)=(r,s).
\]
\noindent Macht man dann weiter
\[
m+n\equiv0\quad(\mathrm{mod.}\ b)\quad\text{und}\quad1+nn'=bc,
\]
\noindent so ist in ganzen Zahlen
\[
Ac=\left(\tfrac{r+n's}{b},\ \tfrac{s-nr}{b}\right)
\]
\noindent u.\,s.\,f.
\clearpage
\SourceLeaf{101}\ArithmeticHead{94}{L}{ARITHMETIK. NACHLASS.}
\noindent\hspace*{\parindent}Allgemein ist nemlich
\[
\begin{aligned}
(p,q)\mathbin{.}(1,m)&=(p+mq,\ q-m'p)\\
&=(p+m'q,\ q-mp)\\
(p,q)\mathbin{.}(P,Q)&=(pP+qQ,\ qP'-pQ')\\
&=(pP-qQ,\ qP'+pQ')\\
&=(pP-qQ,\ q'P+p'Q)\\
&=(pP-q'Q',\ pQ+q'P').
\end{aligned}
\]
\NumberedItem{3}
\noindent\hspace*{\parindent}Steigen wir höher hinauf, so ist alles dieses nur ein sehr spezieller Fall
aus einer viel allgemeineren Untersuchung, die so ausgesprochen werden kann:

\textit{Es sind zu untersuchen die allgemeinen Eigenschaften der unbestimmten Form}
\[
\tfrac12(a+a')xx'+bx'y+b'xy'+\tfrac12(c+c')yy'.
\]
\par\noindent\hspace*{\parindent}Was ist hier der Determinans?\par
\noindent\hspace*{\parindent}Nichts anderes als
\[
bb'-\tfrac14(a+a')(c+c')=\Delta.
\]
\noindent Bezeichnen wir eine solche Form durch
\[
(a,b,c),
\]
\noindent wo $a,c$ reell sind, so geht solche in
\[
\begin{aligned}
&\left(c,mc-b',\ \tfrac{(mc-b)(m'c'-b')-\Delta}{c}\right)\\
={}&(c,mc-b',\ a-mb'-m'b+mm'c)\\
={}&ctt'+(mc-b')t'u+(m'c-b)tu'\\
&\quad +(a-mb'-m'b+mm'c)uu'
\end{aligned}
\]
\noindent über durch die Substitution
\[
\begin{aligned}x&=-u\\y&=t+mu\end{aligned}
\]
\noindent oder umgekehrt
\[
\begin{aligned}t&=mx+y\\u&=-x.\end{aligned}
\]
\clearpage
\SourceLeaf{102}\ArithmeticHead{95}{R}{ZUR THEORIE DER FORMEN.}
\noindent\hspace*{\parindent}Es sind aequivalent
\[
\begin{array}{c|c}
\text{proprie}&\text{improprie}\\
a,b,c&a,-b,c\\
a,ma+b,a-mb'-m'b+mm'c&c,b',a\\
c,-b',a&
\end{array}
\]
\begin{center}\small BEMERKUNG.\end{center}
\begingroup\small
\noindent\hspace*{\parindent}Die vorstehenden Aufzeichnungen stammen aus später Zeit; das geht daraus hervor, daß auf der
Rückseite des Zettels, auf dem sie aufgezeichnet sind, die Bemerkung steht:
\begin{center}\normalsize
\textsc{Jacobi}, Zerlegung der Zahlen in 4 Quadrate,\\
\textsc{Crelle} \textbf{3}, 2; auch \textbf{12}, 2\footnote{Die von \textsc{Gauss} erwähnten Abhandlungen \textsc{Jacobi}s sind: \textit{Note sur la décomposition d’un nombre
donné en quatre carrés}, \textsc{Crelles} Journal f. Mathem. \textbf{3} (1828), S. 191, \textsc{Jacobis} Werke I, S. 245 und \textit{De
compositione numerorum e quattuor quadratis}, \textsc{Crelles} Journal f. Mathem. \textbf{12} (1834), S. 167, \textsc{Jacobis}
Werke VI, S. 245.}.
\end{center}
\noindent\hspace*{\parindent}Die in [2.] erwähnten »Schlüsse von \textsc{Lagrange}« beziehen sich auf das »\textls[120]{Théorème}: Si la somme de
quatre carrés est divisible par un nombre premier plus grand que la racine carrée de la même somme, ce
nombre sera nécessairement égal à la somme de quatre carrés.« (Oeuvres de \textsc{Lagrange} III (1869), S. 193;
der Titel der Abhandlung lautet: \textit{Démonstration d’un théorème d’Arithmétique}).

Bemerkenswert ist in [3.] das Auftreten der heute sogenannten \textsc{Hermite}schen Formen. Vergl. die
bezüglichen Ausführungen in dem Artikel \textbf{12} des \textsc{Bachmann}schen Aufsatzes »Über \textsc{Gauss}’ zahlentheoretische
Arbeiten«.

Die Hilfsformel der Nr. [2.]
\[
(p,q)\mathbin{.}(P,Q)=(pP+qQ,\ qP'-pQ')
\]
\noindent findet sich auch in einer Aufzeichnung des Handbuchs \textbf{19}, Be, S. 147, die Werke III, S. 384 abgedruckt
ist und die sicher aus der Zeit nach 1825 stammt.
\par\hfill\textsc{Schlesinger.}\par
\endgroup
\ClosingRule
\clearpage
